\begin{lemma}
Let \(K\) be a finite field, let \(t\ge2\), and let \(G=G(t,K)\) be the
polarity graph.  Then the star product digraph \(D^*(G)\) is
\(T_{t+1}\)-free.
\end{lemma}
\begin{proof}
By the definition of \(T_{t+1}\)-freeness in
\(\noderef{TransitiveTournamentFree}\), it is enough to rule out an injective
ordered \((t+1)\)-tuple whose earlier vertices send arcs to all later
vertices.  Suppose, for a contradiction, that vertices
\[
  (a_1,b_1),(a_2,b_2),\ldots,(a_{t+1},b_{t+1})
\]
form a copy of \(T_{t+1}\) in \(D^*(G)\) in this order.  Choose nonzero
vectors
\[
  x_1,y_1,\ldots,x_{t+1},y_{t+1}\in K^{t+1}
\]
representing these projective points.  By the definition of
\(\noderef{StarProductDigraph}\) together with the dot-product definition of
the polarity graph in \(\noderef{PolarityGraph}\), for every \(i\) we have
\[
  \langle x_i,y_i\rangle=0,
\]
and for every \(i<j\) we have
\[
  \langle x_i,y_j\rangle=0
  \quad\text{and}\quad
  \langle x_j,y_i\rangle\ne0. \tag{1}
\]
Only the nonzero relations \(\langle x_{i+1},y_i\rangle\ne0\) for
\(1\le i\le t\) will be needed, together with one auxiliary vector.  Since
\(y_{t+1}\ne0\), there is a vector \(x_{t+2}\in K^{t+1}\) with
\(\langle x_{t+2},y_{t+1}\rangle\ne0\).

We claim that \(y_1,\ldots,y_{t+1}\) are linearly independent.  Suppose not.
Let
\[
  \sum_{j=1}^{t+1}\alpha_j y_j=0
\]
be a nontrivial linear relation, and let \(i\) be the least index with
\(\alpha_i\ne0\).  Then
\[
  \sum_{j=i}^{t+1}\alpha_j y_j=0.
\]
Take the inner product with \(x_{i+1}\), where \(x_{t+2}\) is used if
\(i=t+1\).  For every \(j\ge i+1\), relation (1) gives
\(\langle x_{i+1},y_j\rangle=0\); if \(i=t+1\) this is vacuous.  Therefore
\[
  \alpha_i\langle x_{i+1},y_i\rangle=0.
\]
The second relation in (1), or the defining choice of \(x_{t+2}\) when
\(i=t+1\), says that the inner product factor is nonzero.  Hence
\(\alpha_i=0\), contradicting the choice of \(i\).  This proves the claim.

The \(t+1\) vectors \(y_1,\ldots,y_{t+1}\) are therefore a basis of the
\((t+1)\)-dimensional vector space \(K^{t+1}\).  Now \(x_1\) is orthogonal to
every one of them: it is orthogonal to \(y_1\) because \((a_1,b_1)\) is a
vertex of \(D^*(G)\), and it is orthogonal to \(y_j\) for every \(j>1\) by
the arc relation from \((a_1,b_1)\) to \((a_j,b_j)\).  A vector orthogonal to
a basis is the zero vector, so \(x_1=0\).  This contradicts that \(x_1\)
represents a projective point.  Hence no ordered transitive tournament of
size \(t+1\) exists, which is precisely the assertion that \(D^*(G)\) is
\(T_{t+1}\)-free by \(\noderef{TransitiveTournamentFree}\).
\end{proof}
